\documentclass[10pt,a4paper]{amsart}
\usepackage[margin=19mm]{geometry}
\usepackage{amsmath,amssymb,mathtools}
\usepackage[scr=rsfs,cal=euler]{mathalfa}
\usepackage{xfrac}
\usepackage[unicode,hidelinks,bookmarks=false]{hyperref}
\newcommand{\ie}{i.e.\ }
\newcommand{\cf}{cf.\ }
\newcommand{\resp}{respectively\ }
\newcommand{\Subsection}{\subsection}
\providecommand{\nobreakdash}{\nobreak\hbox{-}}
\setlength{\emergencystretch}{2em}
\multlinegap0pt
\usepackage{bm}
\usepackage{bbm}
\usepackage{dsfont}
\usepackage{xspace}
\usepackage{cancel}
\usepackage{mathtools}
\usepackage{amsthm}
\makeatletter
\@ifundefined{thm}{\newtheorem{thm}{Thm}}{}
\@ifundefined{coro}{\newtheorem{coro}[thm]{Corollary}}{}
\@ifundefined{lem}{\newtheorem{lem}[thm]{Lemma}}{}
\makeatother
\DeclareMathOperator{\bl}{\mbox{\sf{bl}}}
\DeclareMathOperator{\fcr}{\mbox{\sf{fcr}}}
\DeclareMathOperator{\idr}{\mbox{\sf{idr}}}
\DeclareMathOperator{\lmax}{\mbox{\sf{lmax}}}
\DeclareMathOperator{\magg}{\mbox{\sf{mag}}}
\DeclareMathOperator{\mcr}{\mbox{\sf{mcr}}}
\DeclareMathOperator{\minn}{\mbox{\sf{min}}}
\DeclareMathOperator{\rmin}{\mbox{\sf{rmin}}}
\DeclareMathOperator{\smc}{\mbox{\sf{smc}}}
\DeclareMathOperator{\ssd}{\mbox{\sf{ssd}}}
\DeclareMathOperator{\zero}{\mbox{\sf{zero}}}
\newcommand{\A}{\mathcal{A}}      
\newcommand{\F}{\mathcal{F}}
\newcommand{\M}{\mathcal{M}}
\newcommand{\PP}{{\mathcal P}}
\newcommand{\V}{\mathcal{V}}
\title[Catalan structures arising from pattern-avoiding Stoimenow m]{Catalan structures arising from pattern-avoiding Stoimenow matchings and other Fishburn objects}
\author{Shuzhen Lv, Sergey Kitaev, Philip B. Zhang}
\date{Source: arXiv:2509.09115v2 (CC BY 4.0)}
\begin{document}
\maketitle

\noindent\emph{Source and licence.} Catalan structures arising from pattern-avoiding Stoimenow matchings and other Fishburn objects. Version arXiv:2509.09115v2 (CC BY 4.0), distributed under the Creative Commons Attribution 4.0 International licence, as recorded in the arXiv licence metadata for this exact version and shown on the abstract page https://arxiv.org/abs/2509.09115v2. Full rights notice: licenses/arxiv-2509.09115-CC-BY-4.0.txt.\par\smallskip

\noindent\textit{Editorial scope.} Counted unit: the Lem whose source label is \texttt{lem-P2-N-101-3142} in the article (source0.tex lines 2561--2563, source label \texttt{lem-P2-N-101-3142}), delivered together with the article's own complete original proof of it (source0.tex lines 2564--2604, one \texttt{proof} environment of 41 lines). No mathematics is written, repaired or re-derived: statement and proof are the article's own text line for line. Required context retained uncounted: coro-P2-irre (lines 1120--1123); sta-P1-case1 (lines 2531--2534). Every definition, display and label of those lines is retained unchanged. Omitted from this package and not redistributed: 1--1119, 1124--2530, 2535--2560, 2605--2793. Nothing inside a retained excerpt is omitted or altered: every retained body is delivered exactly as the article's lines are written, so no omission receipt of any kind accompanies this package. Citations: the retained excerpts carry no citation command at all, so no bibliography entry of the article is transcribed and no reference is redistributed. Reference adaptation: none was needed and none was made; every reference inside the delivered unit resolves inside it, so no unresolved reference or citation is produced. Printed slips kept literal: no printed word, symbol, hypothesis or number of the counted statement or of its proof was altered. Layout and class: the article's own class is \texttt{article} with packages including latexsym, amsthm, amsmath, amssymb, amsfonts, mdframed, graphicx, stmaryrd, gensymb, bm, mathtools, fontspec, multirow, tikz; the wrapper is the lane's pinned \texttt{amsart} editorial wrapper at 10pt on A4, which restates the theorem-like environments the retained text needs and adds the article's own macro definitions that the retained text needs (\texttt{\textbackslash A}, \texttt{\textbackslash F}, \texttt{\textbackslash M}, \texttt{\textbackslash PP}, \texttt{\textbackslash V}, \texttt{\textbackslash bl}, \texttt{\textbackslash fcr}, \texttt{\textbackslash idr}, \texttt{\textbackslash lmax}, \texttt{\textbackslash magg}, \texttt{\textbackslash mcr}, \texttt{\textbackslash minn}, \texttt{\textbackslash rmin}, \texttt{\textbackslash smc}, \texttt{\textbackslash ssd}, \texttt{\textbackslash zero}), with extra wrapper packages (bm, bbm, dsfont, xspace, cancel, mathtools). The delivered numbering is the unit's own. Assets: no retained span contains a figure environment, a graphics-inclusion command, a PDF-page inclusion, a table or a TikZ picture; no image file is copied into this package, the package declares no asset dependency and no image file is redistributed with it. Licence: version arXiv:2509.09115v2 of this article is distributed under the Creative Commons Attribution 4.0 International licence, as recorded in the arXiv licence metadata for this exact version and shown on the abstract page https://arxiv.org/abs/2509.09115v2. A full rights notice is delivered at licenses/arxiv-2509.09115-CC-BY-4.0.txt. Characters: every retained excerpt is pure ASCII, so the delivered engine, XeLaTeX with its default Latin Modern text font, reports no missing character.
	\begin{coro}
		\label{coro-P2-irre}
		We have $|\M^{Ir}_n(P_2)|=C_{n-1}$ for any $n\geq 1$.
	\end{coro}

			\begin{align}\label{sta-P1-case1}
				\mcr(M)= \mcr(M_2)+1, \hspace{0.3cm} \fcr(M)=1, \hspace{0.3cm}\bl(M)= \bl(M_2)+1,\notag \\ 
				\magg(P)=\magg(P_{\text{II}})+1, \hspace{0.3cm} \minn(P)=1, \hspace{0.3cm} \ssd(P)= \ssd(P_{\text{II}})+1.
			\end{align}

	\begin{lem}\label{lem-P2-N-101-3142}
		The statistics $(\mcr, \fcr, \bl)$ over $\M_n(P_2)$, $(\magg(P), \minn, \smc)$ over $\PP_n\bf{(N)}$, $(\lmax, \zero, \rmin)$ over $\A_n(101)$, and $(\rmin, \idr, \lmax)$ over $\F_n(3142)$ have the same distribution.
	\end{lem} 

	\begin{proof} 
		Let $M= \V(M') \oplus M'' \in \M_n(P_2)$, where $M' \in \M_{k-1}(P_2)$ and $M'' \in \M_{n-k}(P_2)$, and $\widetilde{M}=\V(M')$. 
		We also define
		\[
		P=\widetilde{P}\oplus P''\in\PP_{n}\textbf{(N)},\quad
		\alpha=\widetilde{\alpha}\oplus\alpha'' \in \A_n(101),\quad
		\pi=\widetilde{\pi}\oplus\pi'' \in \F_n(3142),
		\]
		to be the respective images of $M$ under the bijections $\Omega$, $\Psi$, and $\Upsilon$. Similarly, $P'$, $\alpha'$, and $\pi'$ denote the corresponding images of $M'$ under these same bijections. We consider the following two cases.
		\begin{itemize}
			\item If $M' =\emptyset$, then $P'=\alpha'=\pi'=\emptyset$. 
			By the decompositions, we have  
			$M=\{[1,2]\}\oplus M''$,  
			$P$ is obtained by adding a new unique minimal element to~$P''$, 
			$\alpha$ by prepending~$0$ to~$\alpha''$, and  
			$\pi=1\oplus \pi''$. 
			Thus, this case follows the same recurrences as~\eqref{sta-P1-case1} for the statistics over $\M_n(P_{2})$, $\PP_{n}\bf{(N)}$, $\A_n(101)$, and $\F_n(3142)$.
			
			\item If $M' \neq \emptyset$, then it is not difficult to see that $\bl(\widetilde{M})= 1$.
			By the construction of~$\V$ in Corollary~\ref{coro-P2-irre}, the newly added reduction arc in~$M'$ yields $\fcr(\widetilde{M})=\fcr(M')+1$
			and leaves the number of maximal crossings unchanged; hence $\mcr(\widetilde{M})=\mcr(M')$.
			Therefore, we have
			\begin{align*}
				\mcr(M)= \mcr(M')+\mcr(M''), \hspace{0.3cm}  \fcr(M)=\fcr(M')+1, \hspace{0.3cm} \bl(M)=\bl(M'')+1.
			\end{align*}
			By the constructions of $\widetilde P$, $\widetilde \alpha$, and $\widetilde \pi$, it is clear that
			$\smc(\widetilde{P})=\rmin(\widetilde{\alpha})=\lmax(\widetilde{\pi})=1$,
			and the $\minn$, $\zero$, and $\idr$ statistics each increase by one in
			$\widetilde{P}$, $\widetilde{\alpha}$ and $\widetilde{\pi}$
			compared to $P'$, $\alpha'$ and $\pi'$, respectively.
			We also see that
			$\magg(\widetilde P)= \magg(P')$,
			$\lmax(\widetilde{\alpha})=\lmax(\alpha')$, and  $\rmin(\widetilde{\pi})=\rmin(\pi')$. 
			Furthermore, by the decompositions of $P$, $\alpha$, and $\pi$, we have
			\begin{align*}
				&\magg(P)=\magg(P')+\magg(P''),\hspace{0.3cm} \minn(P)=\minn(P')+1,\hspace{0.3cm}\smc(P)=\smc(P'')+1,\\ &\lmax(\alpha)=\lmax(\alpha')+\lmax(\alpha''), \hspace{0.3cm}\zero(\alpha)=\zero(\alpha')+1, \hspace{0.3cm} \rmin(\alpha)=\rmin(\alpha'')+1, \\ 
				&\rmin(\pi)=\rmin(\pi')+\rmin(\pi''), \hspace{0.3cm} \idr(\pi)=\idr(\pi')+1, \hspace{0.3cm} \lmax(\pi)=\lmax(\pi'')+1.
			\end{align*}
		\end{itemize}
		The proof now follows by induction on $n$.
	\end{proof}

\end{document}
